\section{Limitations, ethics and broader impact}\label{sec:limits}

\paragraph{Limitations.}
\emph{Scope of the lab.} One narrow task in a sandbox, with fictional hosts, open-weight models from one provider (earlier single-run results with a closed model are not pooled), and swarms of 10 to 100 agents. The numbers describe what these models did here; the \emph{model-dependence} is the finding, not any single rate. \emph{Exposure, not intent.} ``Copied'' means the agent was served the item first; an agent may have used it for a different reason. \emph{One incident.} The ceiling of \bestTop{} comes from one real incident, whose logs happened not to record reads; other swarms will differ, which is why the audit exists. \emph{AI Village is an upper bound.} Room membership is proxied from posting, and the dataset's chat is a channel everyone reads. \emph{Pooled numbers lean on one model.} The pooled benchmark is dominated by the core model; use the per-family results. \emph{Few runs per cell} for most non-core models (two to four), hence wide intervals. \emph{The simulator shares a model family with the tracer} it validates; we test rules outside it, not every rule. \emph{Truth is from our own harness.} The hidden read log is exact for the lab and does not exist for real swarms. \emph{Tokens can be stripped.} We study unintended loss; an agent that deliberately removes or forges tokens defeats a token scheme, and the gateway then relies on closure (Proposition~\ref{prop:gate}). \emph{Gateways need control.} The design assumes the operator serves the shared artifacts and the resource servers; it does not apply to a swarm that reads the open web, and the gate traces access, not ideas.

\paragraph{Ethics, data and safety.}\label{sec:ethics}
All lab runs are offline against reserved \texttt{.invalid} hosts, benign, and against our own simulated world; refusals are logged, and nothing was tested against, probed in, or injected into anyone's live system. Our tools never fetch URLs found in logs. The AI Village dataset is used for research and analysis under its terms (citation of its authors, no training or fine-tuning, no re-identification, notification of publications); we release only aggregates and none of its text. The real incident logs were shared with participants of an event and are marked as not for redistribution; this paper uses aggregates and no handles, page text, URLs or tokens, and publication of those aggregates is \ph{CONFIRMED-WITH-ORGANISERS / DATE}. The audit and the viewer run entirely in the browser or locally, mask strings from a log by default, and share only counts and rates.

\paragraph{Broader impact.}
Tracing how information spreads through agent swarms helps operators find the source of a mistake, re-check only what it reached, and avoid blaming agents that did nothing. The same capability can be misused: logging reads is sensitive (what an agent, or a person behind it, looked at can be private), and per-reader tokens identify readers. We recommend logging reads only for artifacts the operator already controls, retaining registries for a limited time, keeping them separate from the swarm, and not using tracing to profile individuals. Tokens also give an attacker who can read the registry a map of who saw what; the registry should be protected like any access log.
